\section{Related Work, Limitations, and Conclusion}
\label{sec:limits}
\subsection{Positioning}
The earlier diagonal-training design~\cite{mihara2020} motivates rotation-only transposition. We generalize the layout and charge server-side batch maintenance, compact refresh, and optimizer state. ReBoot~\cite{reboot2026} already establishes non-interactive training and local learning; its released implementation is our application baseline, not a source of imported implementation code. HETAL~\cite{hetal2023} studies encrypted transfer learning and efficient matrix products in a different training scope. Our evaluation does not reproduce its workloads or establish a ranking against it.

The cryptographic building blocks are efficient linear transforms and hoisting~\cite{halevi2018}, SSE~\cite{sse2022}, and CKKS~\cite{ckks2017}. SEAL~\cite{seal} provides interoperability checks; OpenFHE~\cite{openfhe2022} supplies the separately executed baseline. HEonGPU~\cite{gpuKeys2025,heongpuBench} is a device-operation reference, not an end-to-end training competitor measured under our protocol. An MIT implementation improves reuse, but licensing and pipeline completeness are not evidence of a cryptographic or learning-performance advance. The approved dependency policy permits MIT, Apache, BSD, and ISC licenses, with one narrow Unicode-3.0 character-table exception for \texttt{unicode-ident} version \texttt{1.0.27}. That component requires its complete Unicode notice; this is not a general Unicode-license exception.

\subsection{Limits and missed goals}
\paragraph{Scope and comparison.}
FULL-1 and FULL-2 supply completed full-setting encrypted runs at two batches with the same initial-weight seed, including their schedules and accumulated error. FULL-2 passes the registered baseline-batch accuracy gate, but misses the common-baseline-bootstrap count target. FULL-1's successful count comparisons are a separate configuration. Full-setting ReBoot CPU wall-clock is unmeasured, and the common-security condition is unmet. Neither these two configurations nor the earlier subsets establish multi-seed encrypted behavior; multi-seed plaintext results are a different evidence class. The observed parameter error does not establish a general error bound. Batch choices trade refresh amortization against learning dynamics. The effective baseline security differs, initialization and bias arithmetic require care, and legacy rows cannot support accuracy claims. ReBoot's in-library layout timing was measured on a different machine and is not a same-CPU baseline; its public-label/client-update path has a weaker confidentiality contract. The official encrypted block-model baselines have not been run: their reported memory demand exceeds the available host budget.

\paragraph{Cost and correctness.}
A key switch at a low level need not cost the same as one at the top level; counts omit memory traffic, cache effects, and key residency. The initial compact-layout assumptions required actual CKKS tests and the first model was optimistic. Replicated copies were once unstable during interactive training; canonical-layout refresh was required to avoid divergence. The earlier paper's published timing total and phase sum disagree, so the continuity demonstration does not silently reconcile them or infer a server-only speedup. The original CPU NTT target and several SEAL-relative key-switch performance targets were missed. The early CUDA NTT spike was not timed against its registered hardware reference; later device-operator measurements do not retroactively supply that measurement. Some SEAL comparisons also have different special-prime/key-switch decompositions. We report these as missed performance targets, not failed numerical gates.

\paragraph{Numerical limits.}
SSE and modulus-band changes were needed because the planned dense-secret overflow range and original prime allocation did not meet the bootstrap precision budget. Unit-scale nonconstant masks were invalid, and exact-scale folding is required to fit the training circuit. The accepted security claim uses specified sieve cost models, not a universal security level; Core-SVP estimates are lower. Dense-secret hybrid search is bounded rather than an exhaustive global optimization. An all-sparse main-key alternative did not meet the selected security criterion and is not used. KAT coverage is limited to the supported SEAL preset, with asymmetric comparison directions; it does not independently validate our bootstrap or larger hybrid preset. Standalone OpenFHE/Lattigo bootstrap precision parity remains absent. Random full-slot CPU precision checks do not prove a negligible failure probability for arbitrary training distributions or activations outside range.

\paragraph{Learning limits.}
The quartic activation's claimed accuracy benefit is unmeasured. The evaluated encrypted trainer omits weight decay and optional softmax/cross-entropy. The separate plaintext paper-setting reconstruction includes weight decay; it is not the faithful encrypted protocol. Bias momentum is not matched to NAG, and faithful bias freezing preserves extra encrypted arithmetic. Fixed-final-epoch Breast Cancer accuracy misses its target; test-selected best epochs are not a remedy. No encrypted Fashion-MNIST or Breast Cancer accuracy is reported. Higher-depth local-block training has plaintext fidelity checks but no completed encrypted comparison. The continuity Iris experiment is interactive and differs in padding, randomization, and normalization; it cannot establish the primary non-interactive claim. Threshold or multi-party FHE, scheme switching, transformer training, and Metal execution are outside the evaluated scope.

\paragraph{Device and implementation limits.}
Full-key GPU training, the matched training parity gate, the RQ3 timing ablation, and GPU bootstrap measurements are pending. Repaired multi-stream latency is measured only for independent products at a capped level; it cannot supply those missing training results. Memory estimates and pool snapshots do not establish a peak. All-off flags do not implement the strongest naive diagonal algorithm. The current training benchmark still uses host bootstrap fallback. Device error handling and numerical replay are useful engineering evidence, but compile-only paths and skipped tests are not physical hardware results. CPU/GPU side-channel resistance, constant-time sampling, secret zeroization, malformed-input hardening, and independent cryptographic audit are outside the current implementation's assurance. Some internal invariants are debug assertions; decode/parameter mismatches can panic. The artifact is research software, not a deployment security claim.

\subsection{Conclusion}
Batched diagonal packing preserves a rotation-only transpose while supporting compact non-interactive optimizer updates. The completed full-setting runs show a larger-batch count advantage and baseline-batch accuracy within tolerance, together with a missed normalized count target at the baseline batch. The security mismatch and unmeasured full-setting baseline time prevent an unconditional superiority claim. The device microbenchmark also shows that stream and pool choices matter, without resolving training performance. The resulting evidence supports a conditional re-evaluation of the layout and identifies the matched-security, baseline-timing, and hardware experiments still needed.
